% Options for packages loaded elsewhere
\PassOptionsToPackage{unicode}{hyperref}
\PassOptionsToPackage{hyphens}{url}
\documentclass[
]{article}
\usepackage{xcolor}
\usepackage[margin=1in]{geometry}
\usepackage{amsmath,amssymb}
\setcounter{secnumdepth}{-\maxdimen} % remove section numbering
\usepackage{iftex}
\ifPDFTeX
  \usepackage[T1]{fontenc}
  \usepackage[utf8]{inputenc}
  \usepackage{textcomp} % provide euro and other symbols
\else % if luatex or xetex
  \usepackage{unicode-math} % this also loads fontspec
  \defaultfontfeatures{Scale=MatchLowercase}
  \defaultfontfeatures[\rmfamily]{Ligatures=TeX,Scale=1}
\fi
\usepackage{lmodern}
\ifPDFTeX\else
  % xetex/luatex font selection
\fi
% Use upquote if available, for straight quotes in verbatim environments
\IfFileExists{upquote.sty}{\usepackage{upquote}}{}
\IfFileExists{microtype.sty}{% use microtype if available
  \usepackage[]{microtype}
  \UseMicrotypeSet[protrusion]{basicmath} % disable protrusion for tt fonts
}{}
\makeatletter
\@ifundefined{KOMAClassName}{% if non-KOMA class
  \IfFileExists{parskip.sty}{%
    \usepackage{parskip}
  }{% else
    \setlength{\parindent}{0pt}
    \setlength{\parskip}{6pt plus 2pt minus 1pt}}
}{% if KOMA class
  \KOMAoptions{parskip=half}}
\makeatother
\usepackage{color}
\usepackage{fancyvrb}
\newcommand{\VerbBar}{|}
\newcommand{\VERB}{\Verb[commandchars=\\\{\}]}
\DefineVerbatimEnvironment{Highlighting}{Verbatim}{commandchars=\\\{\}}
% Add ',fontsize=\small' for more characters per line
\usepackage{framed}
\definecolor{shadecolor}{RGB}{248,248,248}
\newenvironment{Shaded}{\begin{snugshade}}{\end{snugshade}}
\newcommand{\AlertTok}[1]{\textcolor[rgb]{0.94,0.16,0.16}{#1}}
\newcommand{\AnnotationTok}[1]{\textcolor[rgb]{0.56,0.35,0.01}{\textbf{\textit{#1}}}}
\newcommand{\AttributeTok}[1]{\textcolor[rgb]{0.13,0.29,0.53}{#1}}
\newcommand{\BaseNTok}[1]{\textcolor[rgb]{0.00,0.00,0.81}{#1}}
\newcommand{\BuiltInTok}[1]{#1}
\newcommand{\CharTok}[1]{\textcolor[rgb]{0.31,0.60,0.02}{#1}}
\newcommand{\CommentTok}[1]{\textcolor[rgb]{0.56,0.35,0.01}{\textit{#1}}}
\newcommand{\CommentVarTok}[1]{\textcolor[rgb]{0.56,0.35,0.01}{\textbf{\textit{#1}}}}
\newcommand{\ConstantTok}[1]{\textcolor[rgb]{0.56,0.35,0.01}{#1}}
\newcommand{\ControlFlowTok}[1]{\textcolor[rgb]{0.13,0.29,0.53}{\textbf{#1}}}
\newcommand{\DataTypeTok}[1]{\textcolor[rgb]{0.13,0.29,0.53}{#1}}
\newcommand{\DecValTok}[1]{\textcolor[rgb]{0.00,0.00,0.81}{#1}}
\newcommand{\DocumentationTok}[1]{\textcolor[rgb]{0.56,0.35,0.01}{\textbf{\textit{#1}}}}
\newcommand{\ErrorTok}[1]{\textcolor[rgb]{0.64,0.00,0.00}{\textbf{#1}}}
\newcommand{\ExtensionTok}[1]{#1}
\newcommand{\FloatTok}[1]{\textcolor[rgb]{0.00,0.00,0.81}{#1}}
\newcommand{\FunctionTok}[1]{\textcolor[rgb]{0.13,0.29,0.53}{\textbf{#1}}}
\newcommand{\ImportTok}[1]{#1}
\newcommand{\InformationTok}[1]{\textcolor[rgb]{0.56,0.35,0.01}{\textbf{\textit{#1}}}}
\newcommand{\KeywordTok}[1]{\textcolor[rgb]{0.13,0.29,0.53}{\textbf{#1}}}
\newcommand{\NormalTok}[1]{#1}
\newcommand{\OperatorTok}[1]{\textcolor[rgb]{0.81,0.36,0.00}{\textbf{#1}}}
\newcommand{\OtherTok}[1]{\textcolor[rgb]{0.56,0.35,0.01}{#1}}
\newcommand{\PreprocessorTok}[1]{\textcolor[rgb]{0.56,0.35,0.01}{\textit{#1}}}
\newcommand{\RegionMarkerTok}[1]{#1}
\newcommand{\SpecialCharTok}[1]{\textcolor[rgb]{0.81,0.36,0.00}{\textbf{#1}}}
\newcommand{\SpecialStringTok}[1]{\textcolor[rgb]{0.31,0.60,0.02}{#1}}
\newcommand{\StringTok}[1]{\textcolor[rgb]{0.31,0.60,0.02}{#1}}
\newcommand{\VariableTok}[1]{\textcolor[rgb]{0.00,0.00,0.00}{#1}}
\newcommand{\VerbatimStringTok}[1]{\textcolor[rgb]{0.31,0.60,0.02}{#1}}
\newcommand{\WarningTok}[1]{\textcolor[rgb]{0.56,0.35,0.01}{\textbf{\textit{#1}}}}
\usepackage{graphicx}
\makeatletter
\newsavebox\pandoc@box
\newcommand*\pandocbounded[1]{% scales image to fit in text height/width
  \sbox\pandoc@box{#1}%
  \Gscale@div\@tempa{\textheight}{\dimexpr\ht\pandoc@box+\dp\pandoc@box\relax}%
  \Gscale@div\@tempb{\linewidth}{\wd\pandoc@box}%
  \ifdim\@tempb\p@<\@tempa\p@\let\@tempa\@tempb\fi% select the smaller of both
  \ifdim\@tempa\p@<\p@\scalebox{\@tempa}{\usebox\pandoc@box}%
  \else\usebox{\pandoc@box}%
  \fi%
}
% Set default figure placement to htbp
\def\fps@figure{htbp}
\makeatother
\setlength{\emergencystretch}{3em} % prevent overfull lines
\providecommand{\tightlist}{%
  \setlength{\itemsep}{0pt}\setlength{\parskip}{0pt}}
\usepackage{bookmark}
\IfFileExists{xurl.sty}{\usepackage{xurl}}{} % add URL line breaks if available
\urlstyle{same}
\hypersetup{
  pdftitle={Maximum likelihood estimation},
  hidelinks,
  pdfcreator={LaTeX via pandoc}}

\title{Maximum likelihood estimation}
\author{}
\date{\vspace{-2.5em}}

\begin{document}
\maketitle

\textbf{Inference labs} use the five-step template: Hypothesis →
Visualise → Assumptions → Conduct → Conclude.

\subsection{Learning objectives}\label{learning-objectives}

\begin{itemize}
\tightlist
\item
  Write a log-likelihood for a simple model and find its maximum.
\item
  Use Fisher information to obtain an asymptotic standard error and
  confidence interval.
\item
  Relate the MLE to intuitive estimators (sample mean, sample rate).
\end{itemize}

\subsection{Prerequisites}\label{prerequisites}

Lab 3.1.

\subsection{Background}\label{background}

Maximum likelihood is the workhorse of parametric inference. Given a
parametric family for the data --- normal, Poisson, binomial --- the MLE
is the parameter value that makes the observed data most probable. Under
mild regularity conditions the MLE is consistent (it converges to the
truth), asymptotically unbiased, and asymptotically normal with variance
given by the inverse of Fisher information.

Fisher information quantifies how sharply the log-likelihood peaks at
its maximum. A sharply peaked likelihood gives a precise estimate (small
variance); a flat one gives a vague estimate. The standard error of an
MLE is the square root of 1/\emph{n} times the inverse of the Fisher
information per observation.

In the simple cases we treat here the MLE coincides with the intuitive
estimator --- the mean for a normal, the sample rate for a Poisson ---
but the machinery generalises to any likelihood.

\subsection{Setup}\label{setup}

\begin{Shaded}
\begin{Highlighting}[]
\FunctionTok{library}\NormalTok{(tidyverse)}
\FunctionTok{set.seed}\NormalTok{(}\DecValTok{42}\NormalTok{)}
\FunctionTok{theme\_set}\NormalTok{(}\FunctionTok{theme\_minimal}\NormalTok{(}\AttributeTok{base\_size =} \DecValTok{12}\NormalTok{))}
\end{Highlighting}
\end{Shaded}

\subsection{1. Hypothesis}\label{hypothesis}

Estimate (i) the mean of a normal and (ii) the rate of a Poisson by
maximum likelihood, recover standard errors from Fisher information, and
build Wald confidence intervals.

\subsection{2. Visualise}\label{visualise}

Normal data with known σ = 2.

\begin{Shaded}
\begin{Highlighting}[]
\NormalTok{mu\_true }\OtherTok{\textless{}{-}} \DecValTok{5}
\NormalTok{sigma }\OtherTok{\textless{}{-}} \DecValTok{2}
\NormalTok{x }\OtherTok{\textless{}{-}} \FunctionTok{rnorm}\NormalTok{(n, }\AttributeTok{mean =}\NormalTok{ mu\_true, }\AttributeTok{sd =}\NormalTok{ sigma)}
\end{Highlighting}
\end{Shaded}

\begin{Shaded}
\begin{Highlighting}[]
\NormalTok{grid }\OtherTok{\textless{}{-}} \FunctionTok{tibble}\NormalTok{(}\AttributeTok{mu =} \FunctionTok{seq}\NormalTok{(}\FloatTok{3.5}\NormalTok{, }\FloatTok{6.5}\NormalTok{, }\AttributeTok{length.out =} \DecValTok{300}\NormalTok{)) }\SpecialCharTok{|\textgreater{}}
  \FunctionTok{mutate}\NormalTok{(}\AttributeTok{loglik =} \FunctionTok{sapply}\NormalTok{(mu, }\ControlFlowTok{function}\NormalTok{(m) }\FunctionTok{sum}\NormalTok{(}\FunctionTok{dnorm}\NormalTok{(x, m, sigma, }\AttributeTok{log =} \ConstantTok{TRUE}\NormalTok{))))}
\FunctionTok{ggplot}\NormalTok{(grid, }\FunctionTok{aes}\NormalTok{(mu, loglik)) }\SpecialCharTok{+}
  \FunctionTok{geom\_line}\NormalTok{() }\SpecialCharTok{+}
  \FunctionTok{geom\_vline}\NormalTok{(}\AttributeTok{xintercept =} \FunctionTok{mean}\NormalTok{(x), }\AttributeTok{linetype =} \DecValTok{2}\NormalTok{, }\AttributeTok{colour =} \StringTok{"firebrick"}\NormalTok{) }\SpecialCharTok{+}
  \FunctionTok{labs}\NormalTok{(}\AttributeTok{x =} \FunctionTok{expression}\NormalTok{(mu), }\AttributeTok{y =} \StringTok{"log{-}likelihood"}\NormalTok{)}
\end{Highlighting}
\end{Shaded}

The likelihood is peaked at the sample mean --- that is the MLE.

\subsection{3. Assumptions}\label{assumptions}

Normal model: iid, constant σ. Poisson model: iid counts, constant rate
over a fixed exposure. Fisher information is evaluated \emph{at the MLE}
for the Wald standard error.

\subsection{4. Conduct}\label{conduct}

\subsubsection{Normal mean by MLE}\label{normal-mean-by-mle}

\begin{Shaded}
\begin{Highlighting}[]
\CommentTok{\# Fisher information for mu in normal with known sigma is n / sigma\^{}2.}
\NormalTok{I\_mu }\OtherTok{\textless{}{-}}\NormalTok{ n }\SpecialCharTok{/}\NormalTok{ sigma}\SpecialCharTok{\^{}}\DecValTok{2}
\NormalTok{se\_mu }\OtherTok{\textless{}{-}} \DecValTok{1} \SpecialCharTok{/} \FunctionTok{sqrt}\NormalTok{(I\_mu)}
\NormalTok{ci\_mu }\OtherTok{\textless{}{-}}\NormalTok{ mle\_mu }\SpecialCharTok{+} \FunctionTok{c}\NormalTok{(}\SpecialCharTok{{-}}\DecValTok{1}\NormalTok{, }\DecValTok{1}\NormalTok{) }\SpecialCharTok{*} \FunctionTok{qnorm}\NormalTok{(}\FloatTok{0.975}\NormalTok{) }\SpecialCharTok{*}\NormalTok{ se\_mu}
\FunctionTok{tibble}\NormalTok{(}\AttributeTok{mle =}\NormalTok{ mle\_mu, }\AttributeTok{se =}\NormalTok{ se\_mu,}
       \AttributeTok{ci\_low =}\NormalTok{ ci\_mu[}\DecValTok{1}\NormalTok{], }\AttributeTok{ci\_high =}\NormalTok{ ci\_mu[}\DecValTok{2}\NormalTok{])}
\end{Highlighting}
\end{Shaded}

\subsubsection{Normal mean by numeric optimisation (sanity
check)}\label{normal-mean-by-numeric-optimisation-sanity-check}

\begin{Shaded}
\begin{Highlighting}[]
\NormalTok{opt }\OtherTok{\textless{}{-}} \FunctionTok{optim}\NormalTok{(}\DecValTok{0}\NormalTok{, nll, }\AttributeTok{method =} \StringTok{"Brent"}\NormalTok{, }\AttributeTok{lower =} \DecValTok{0}\NormalTok{, }\AttributeTok{upper =} \DecValTok{10}\NormalTok{, }\AttributeTok{hessian =} \ConstantTok{TRUE}\NormalTok{)}
\NormalTok{opt}\SpecialCharTok{$}\NormalTok{par        }\CommentTok{\# MLE}
\FunctionTok{sqrt}\NormalTok{(}\DecValTok{1} \SpecialCharTok{/}\NormalTok{ opt}\SpecialCharTok{$}\NormalTok{hessian[}\DecValTok{1}\NormalTok{, }\DecValTok{1}\NormalTok{])  }\CommentTok{\# SE from observed information}
\end{Highlighting}
\end{Shaded}

\subsubsection{Poisson rate by MLE}\label{poisson-rate-by-mle}

\begin{Shaded}
\begin{Highlighting}[]
\NormalTok{mle\_lambda }\OtherTok{\textless{}{-}} \FunctionTok{mean}\NormalTok{(y)}
\CommentTok{\# Fisher information for lambda: n / lambda. At MLE use 1/mle\_lambda per obs.}
\NormalTok{I\_lambda }\OtherTok{\textless{}{-}} \FunctionTok{length}\NormalTok{(y) }\SpecialCharTok{/}\NormalTok{ mle\_lambda}
\NormalTok{se\_lambda }\OtherTok{\textless{}{-}} \DecValTok{1} \SpecialCharTok{/} \FunctionTok{sqrt}\NormalTok{(I\_lambda)}
\NormalTok{ci\_lambda }\OtherTok{\textless{}{-}}\NormalTok{ mle\_lambda }\SpecialCharTok{+} \FunctionTok{c}\NormalTok{(}\SpecialCharTok{{-}}\DecValTok{1}\NormalTok{, }\DecValTok{1}\NormalTok{) }\SpecialCharTok{*} \FunctionTok{qnorm}\NormalTok{(}\FloatTok{0.975}\NormalTok{) }\SpecialCharTok{*}\NormalTok{ se\_lambda}
\FunctionTok{tibble}\NormalTok{(}\AttributeTok{mle =}\NormalTok{ mle\_lambda, }\AttributeTok{se =}\NormalTok{ se\_lambda,}
       \AttributeTok{ci\_low =}\NormalTok{ ci\_lambda[}\DecValTok{1}\NormalTok{], }\AttributeTok{ci\_high =}\NormalTok{ ci\_lambda[}\DecValTok{2}\NormalTok{])}
\end{Highlighting}
\end{Shaded}

\subsubsection{Likelihood surface for a two-parameter model (μ,
σ)}\label{likelihood-surface-for-a-two-parameter-model-ux3bc-ux3c3}

\begin{Shaded}
\begin{Highlighting}[]
\NormalTok{  mu }\OtherTok{=} \FunctionTok{seq}\NormalTok{(}\FloatTok{3.5}\NormalTok{, }\FloatTok{6.5}\NormalTok{, }\AttributeTok{length.out =} \DecValTok{80}\NormalTok{),}
\NormalTok{  s  }\OtherTok{=} \FunctionTok{seq}\NormalTok{(}\DecValTok{1}\NormalTok{, }\FloatTok{3.5}\NormalTok{, }\AttributeTok{length.out =} \DecValTok{80}\NormalTok{)}
\ErrorTok{)} \SpecialCharTok{|\textgreater{}}
  \FunctionTok{rowwise}\NormalTok{() }\SpecialCharTok{|\textgreater{}}
  \FunctionTok{mutate}\NormalTok{(}\AttributeTok{ll =} \FunctionTok{sum}\NormalTok{(}\FunctionTok{dnorm}\NormalTok{(x, mu, s, }\AttributeTok{log =} \ConstantTok{TRUE}\NormalTok{))) }\SpecialCharTok{|\textgreater{}}
  \FunctionTok{ungroup}\NormalTok{()}
\end{Highlighting}
\end{Shaded}

\begin{Shaded}
\begin{Highlighting}[]
\FunctionTok{ggplot}\NormalTok{(surf, }\FunctionTok{aes}\NormalTok{(mu, s, }\AttributeTok{fill =}\NormalTok{ ll)) }\SpecialCharTok{+}
  \FunctionTok{geom\_tile}\NormalTok{() }\SpecialCharTok{+}
  \FunctionTok{geom\_contour}\NormalTok{(}\FunctionTok{aes}\NormalTok{(}\AttributeTok{z =}\NormalTok{ ll), }\AttributeTok{colour =} \StringTok{"white"}\NormalTok{) }\SpecialCharTok{+}
  \FunctionTok{scale\_fill\_viridis\_c}\NormalTok{() }\SpecialCharTok{+}
  \FunctionTok{labs}\NormalTok{(}\AttributeTok{x =} \FunctionTok{expression}\NormalTok{(mu), }\AttributeTok{y =} \FunctionTok{expression}\NormalTok{(sigma), }\AttributeTok{fill =} \StringTok{"log{-}L"}\NormalTok{)}
\end{Highlighting}
\end{Shaded}

\subsection{5. Concluding statement}\label{concluding-statement}

\begin{quote}
The MLE of the normal mean from \emph{n} = \texttt{n} observations at σ
= 2 was \texttt{round(mle\_mu,\ 3)} (Fisher-information SE =
\texttt{round(se\_mu,\ 3)}; 95\% Wald CI
\texttt{round(ci\_mu{[}1{]},\ 2)} to \texttt{round(ci\_mu{[}2{]},\ 2)}),
recovering the true μ = 5. The MLE of the Poisson rate from \emph{n} =
100 observations was \texttt{round(mle\_lambda,\ 3)} (SE
\texttt{round(se\_lambda,\ 3)}; CI \texttt{round(ci\_lambda{[}1{]},\ 2)}
to \texttt{round(ci\_lambda{[}2{]},\ 2)}). Observed-information SEs from
\texttt{optim} agreed with the expected-information calculation.
\end{quote}

The MLE is the general-purpose engine behind linear regression, logistic
regression, Cox models, and the bulk of parametric inference in this
course. The machinery --- likelihood, information, SE --- is the same in
each case.

\subsection{Common pitfalls}\label{common-pitfalls}

\begin{itemize}
\tightlist
\item
  Reporting a Wald CI when the sample is small. Profile CIs are more
  accurate in that regime.
\item
  Evaluating Fisher information at a wrong value of the parameter. Use
  the MLE.
\item
  Forgetting that the MLE is not always unbiased for finite \emph{n} ---
  only asymptotically.
\end{itemize}

\subsection{Further reading}\label{further-reading}

\begin{itemize}
\tightlist
\item
  Casella G, Berger RL. \emph{Statistical Inference}, chapter on point
  estimation.
\item
  Cox DR. \emph{Principles of Statistical Inference}.
\end{itemize}

\subsection{Session info}\label{session-info}

\begin{Shaded}
\begin{Highlighting}[]

\end{Highlighting}
\end{Shaded}

\subsection{See also --- chapter index}\label{see-also-chapter-index}

\begin{itemize}
\tightlist
\item
  \href{../index.Rmd}{Methods in this chapter}
\item
  \href{../../../ch00-find-your-method.Rmd}{Find your method (Chapter
  0)}
\end{itemize}

\end{document}
